\section{基于Bagging的集成}
\label{sec:ensemble-bagging}

决策树的上层分裂可能随少量训练记录变化而改变，预测也随之波动。对同一总体的多份训练集分别建模并平均，可以削弱其中不共享的波动；实际通常只有一份数据，因而需要近似这种重复训练。\term{装袋法}（bootstrap aggregating, Bagging）用自助采样从现有训练集构造多份样本，分别训练成员，再用固定规则聚合。重采样使训练波动转为成员差异，固定聚合再削弱其中部分变化；成员共享的表示偏差则继续保留\scite{breiman1996bagging}。

\subsection{Bagging}
\label{sec:ensemble-bagging-bootstrap}

\subsubsection{自助样本与聚合预测}

对第$m$个成员，从$\{1,\ldots,n\}$中独立、等概率、有放回地抽取$n$次索引，所得多重集合记为$I^{(m)}$。有放回抽样使部分记录重复进入成员训练，也使部分记录留在该成员之外。学习算法$\mathcal{A}$在相应自助样本上训练出
\[
  h_m=\mathcal{A}\!\left(D_n^{(m)}\right),
  \qquad
  D_n^{(m)}=\{(\bm{x}_i,y_i):i\in I^{(m)}\}.
\]
回归通常平均$h_m(\bm{x})$；分类可以投票，也可以在成员输出具有同一语义时平均逐类得分或概率。可以改变自助样本大小或采用其他成员算法。此时仍需检验重采样是否使成员在新样本上产生互补预测。

一份大小为$n$的标准自助样本没有抽中原始样本$i$的概率为
\begin{equation}
  \mathbb{P}\!\left(i\notin I^{(m)}\right)
  =
  \left(1-\frac{1}{n}\right)^n
  \longrightarrow e^{-1}.
  \label{eq:ensemble-oob-probability}
\end{equation}
因此，每个成员平均约有$36.8\%$的原始样本未参与其训练。这些样本构成该成员的\term{袋外样本}（out-of-bag sample, OOB sample）。反过来，对每个训练样本$i$，可以只收集没有使用它训练的成员预测。令
\[
  \mathcal{M}_{i}^{\mathrm{OOB}}
  =\{m:i\notin I^{(m)}\},
\]
当$\mathcal{M}_{i}^{\mathrm{OOB}}$非空时，回归的袋外预测可写为
\begin{equation}
  \widehat{y}_{i}^{\mathrm{OOB}}
  =
  \frac{1}{|\mathcal{M}_{i}^{\mathrm{OOB}}|}
  \sum_{m\in\mathcal{M}_{i}^{\mathrm{OOB}}}h_m(\bm{x}_i),
  \label{eq:ensemble-oob-prediction}
\end{equation}
分类则对同一成员集合投票或平均类别输出。图\ref{fig:ensemble-bagging-oob}按样本收集符合条件的成员：每个样本的袋外预测由其未参与训练的成员子集给出，子集随样本变化。

\begin{figure*}[htbp]
  \centering
  \includegraphics[width=130.048mm]{figures/03-models/01-classic-models/generated/ensemble-bagging-oob/figure.png}
  \caption{Bagging训练与袋外预测。每个成员只使用其自助样本训练；对训练样本$\bm{x}_i$，袋外预测只聚合未抽中$i$的成员。索引与集合展示样本和成员的对应关系；成员数有限时，各样本得到的袋外预测数量不同。}
  \label{fig:ensemble-bagging-oob}
\end{figure*}

\subsubsection{成员数量与误差来源}

给定原数据后，独立的重采样及训练可分配给不同计算单元。评价并行收益时，需要同时计入总工作量与实际经过时间：若单个成员的训练和预测成本分别为$C_{\mathrm{train}}$与$C_{\mathrm{pred}}$，朴素实现的总工作量仍约随$M C_{\mathrm{train}}$和$M C_{\mathrm{pred}}$增长。模型存储也通常随成员数量增加。

在固定训练集上增加独立随机成员，主要减少有限次抽样的Monte Carlo误差；等权平均逐渐逼近该随机训练机制的期望预测。原数据的抽样误差和模型失配仍影响这一期望，共同误差也会留在组合中。实践中可以观察袋外损失随成员数的变化，在曲线趋于稳定后停止增加成员，但这个停止规则也属于模型选择过程。

单树复杂度控制每份样本上的拟合，树数控制随机预测的平均程度。较深树可以表达更细的分区，也可能对扰动更敏感；叶规模与深度限制则可减少极小叶和噪声标签的影响。叶最小样本量、最大深度、类别权重和自助样本大小仍须与成员数共同选择，尤其在标签噪声、极小叶或类别极不平衡时。

\subsubsection{袋外预测的用途与边界}

袋外预测使当前标签与产生该预测的成员训练分离，可用于绘制树数—损失曲线、内部误差估计和置换重要性。可用成员数随机变化，自助训练集又相互重叠，这一数据结构与固定折的交叉验证不同。成员数很小时，某些样本可能没有足够的袋外预测。

反复用袋外结果选择结构、树数或特征处理，会使袋外误差成为开发证据；最终冻结流程需要另用测试数据评价。时间序列、分组样本或空间相关数据还要求自助抽样尊重依赖单位；逐行重采样会破坏原有数据结构。袋外结果的解释还依赖抽样单位与未来使用的数据制度。

\subsection{随机森林}
\label{sec:ensemble-random-forest}

树Bagging改变样本，却可能保留相同的强特征和上层分裂；这些共同选择使成员预测继续相关。\term{随机森林}（random forest）在树的自助采样之外加入节点级特征随机化：每个节点先随机抽取一组候选特征，再只在这组特征中搜索分裂。树仍按第\ref{sec:regression-trees}节或第\ref{sec:classification-trees}节的任务准则形成叶输出，改变的是每个成员可参与分裂搜索的信息\scite{breiman2001randomforests}。

节点候选特征数同时影响单树可用的信息和不同树的分裂差异。候选数接近全部特征时，森林更接近普通的树Bagging；候选数较小时，不同树更可能采用不同分裂，但单棵树也更容易错过当前节点上的有效特征。合适的候选规模取决于特征冗余、样本量和树深，可在固定验证协议下比较。

固定训练集并增加树数，可以减少森林平均预测的随机波动；所得极限仍由当前样本及树构造决定。总体风险另外受分布与估计条件控制，更多树也继续消耗存储和推理计算。随机森林的一致性结论通常针对特定的树构造、抽样比例与分布条件；实际软件在缺失值、类别特征、停止规则和并行实现上的差异，可能不满足同一组理论假设\scite{scornet2015consistency}。

\subsubsection{特征重要性的解释}

树集成常附带两类特征重要性。\term{不纯度下降重要性}（impurity-decrease importance）累计某个特征参与分裂时带来的加权目标下降，计算方便，却可能偏向候选切分较多的特征。相关变量可以在分裂中相互替代，累计下降反映实际选中的路径，因而与变量的独立贡献不同。

\term{置换重要性}（permutation importance）固定已训练模型，在评价数据上打乱一个特征，再观察预测损失增加多少。这个量直接对应已拟合模型在所选扰动下对该特征的依赖，但打乱会破坏该特征与其他特征的联合关系。高度相关的特征可能互相替代，使单个特征的置换影响看起来很小；不现实的特征组合也可能夸大损失变化。这些量分别描述训练分裂和预测扰动中的模型依赖，解释时需说明评价数据、损失及置换方式；因果判断另需因果假设与证据\scite{strobl2007bias}。

图\ref{fig:ensemble-tree-randomization}沿一棵树的构造路径比较三种随机化位置。树Bagging主要改变成员看到的训练样本；随机森林还在每个节点抽取候选特征；Extra Trees继续把随机性放入候选阈值。三者继续以任务准则比较分裂，差别在于进入贪心搜索的样本、特征和阈值。

\begin{figure*}[htbp]
  \centering
  \includegraphics[width=150.199mm]{figures/03-models/01-classic-models/generated/ensemble-tree-randomization/figure.png}
  \caption{树Bagging、随机森林与Extra Trees的随机化位置。每栏从样本来源出发，依次展示候选特征、候选阈值与成员树；候选选择中的随机性和任务准则约束分别标在相应步骤。每栏首框直接标出训练样本来源：前两栏采用自助采样，Extra Trees栏采用原始方法的完整训练集设定；具体实现也可能额外启用自助采样。}
  \label{fig:ensemble-tree-randomization}
\end{figure*}

\subsection{极度随机树（Extra Trees）}
\label{sec:ensemble-extra-trees}

随机森林已随机选择候选特征，阈值仍通过当前节点的数据搜索。\term{极度随机树}（extremely randomized trees, Extra Trees）把随机性进一步放入阈值选择：在每个节点，对随机选出的候选特征生成随机切分点，再从这些随机候选中选择得分较好的切分。原始方法通常让各棵树使用完整训练集，把成员差异主要交给特征与阈值随机化；具体软件也可能允许额外使用自助采样\scite{geurts2006extra}。

随机阈值减少了穷举候选切分的工作，也会让单棵树偏离当前样本上的最佳贪心分裂。组合是否改善，取决于单树预测能力下降与成员依赖减弱之间的平衡。候选特征数、每个特征的随机阈值数、叶最小样本量以及是否采用自助采样共同决定这个平衡；增加随机化能否降低测试损失，需要在同一数据与预算下检验。

这些扩展可从“扰动发生在哪里”比较：在本节的宽泛机制分类中，样本、特征或阈值扰动都用于产生可独立训练的成员，再按固定规则聚合。标准Bagging则特指自助采样聚合。随机子空间只扰动成员可见的特征，Pasting以无放回子样本替代自助样本，随机权重则让同一记录在不同成员中具有不同影响。改变扰动方式后，袋外预测的可用条件也随是否采用自助采样而变化。

\subsection{模型对比与总结}
\label{sec:ensemble-bagging-comparison}

表\ref{tab:ensemble-bagging-comparison}固定树为成员，比较三种方法怎样产生差异。普通Bagging只改变训练样本；随机森林再限制节点可见的候选特征；Extra Trees进一步随机化候选阈值。三种方式还可能采用不同抽样方案，单树误差和搜索成本需与成员依赖一起比较。

\begin{table*}[htbp]
  \centering
  \small
  \begin{tabularx}{\linewidth}{@{}>{\raggedright\arraybackslash}p{24mm}*{6}{>{\raggedright\arraybackslash}X}@{}}
    \toprule
    方法 & 样本扰动 & 节点候选特征 & 阈值选择 & 袋外预测 & 成员训练 & 主要边界 \\
    \midrule
    树Bagging & 通常自助采样 & 全部特征 & 在候选中搜索 & 标准自助设定下可用 & 成员间可并行 & 强特征可能使树高度相关 \\
    随机森林 & 通常自助采样 & 随机特征子集 & 在子集内搜索 & 标准自助设定下可用 & 成员间可并行 & 候选数过小会削弱单树 \\
    Extra Trees & 原始方法使用完整训练集；实现可选自助采样 & 随机特征子集 & 生成随机阈值后择优 & 仅在采用自助采样时可用 & 成员间可并行 & 过强随机化会增大单树误差 \\
    \bottomrule
  \end{tabularx}
  \caption{以树为成员时，Bagging、随机森林与Extra Trees的机制比较。表中列出基本构造；比较具体实现时，还需记录抽样设置，并固定树复杂度、成员数和计算预算。}
  \label{tab:ensemble-bagging-comparison}
\end{table*}

Bagging先产生训练差异，再聚合预测以削弱部分波动。固定聚合保留成员共享的遗漏；若希望让新成员针对这些不足训练，就需要根据当前组合产生修正信号。下一节据此转向序贯提升。
